% ============================================================================
% ⛔ RETIRED 2026-06-23 — DO NOT SUBMIT. Archived for record only.
%
% DOUBLE KILL:
%  (1) The proof below (ultraproduct of Diracs on P(N)) is FATALLY BROKEN — the
%      same diagonal-copy / aleph_1-saturation defect for which this note was
%      previously WITHDRAWN from PAMS. The witness A_k={k,k+1,...} does NOT
%      decrease to 0 in the ultraproduct's Boolean algebra (it is trapped above
%      [d], d(n)=A_n, with mu*([d])=1); and by aleph_1-saturation the ultraproduct
%      charge IS sigma-additive (Loeb phenomenon), so NO ultraproduct over a
%      nonprincipal ultrafilter on N can witness the result.
%  (2) A CORRECT proof exists (ultrapower of the ultrafilter charge: M=(P(N),l)
%      non-sigma-additive, N=M^N/U saturated hence sigma-additive, M elem.equiv. N
%      => not FO-axiomatizable) but the RESULT IS FOLKLORE: stated in
%      Fagin-Halpern-Megiddo 1990 (this note's own cited reference) and proved via
%      Loeb 1975 saturation; standard motivation across model-theory-of-measure.
%      Clears no contribution bar (Type 1 + Type 2 FAIL).
%
% The CE non-derivability FACT survives only as internal synthesis in Paper I
% ("Probability from Observation"), which cites FHM 1990 / Los directly.
% Full verdict: memory countable_additivity_note_proof_broken.md.
% ============================================================================
\documentclass[11pt]{article}

% ---------- Packages ----------
\usepackage[a4paper,margin=1in]{geometry}
\usepackage{amsmath,amssymb,amsthm,mathtools}
\usepackage{hyperref}
\usepackage{enumitem}
\usepackage[numbers]{natbib}

% ---------- Theorem environments ----------
\newtheorem{proposition}{Proposition}
\newtheorem{corollary}[proposition]{Corollary}

\theoremstyle{definition}
\newtheorem{definition}[proposition]{Definition}
\newtheorem{remark}[proposition]{Remark}

% ---------- Commands ----------
\newcommand{\Lmu}{\mathcal{L}_{\mathrm{BA},\mu}}
\newcommand{\Tfa}{T_{\mathrm{fa}}}
\newcommand{\N}{\mathbb{N}}

% ---------- Title ----------
\title{Countable Additivity is Not First-Order Axiomatizable}
\author{Patrick S.\ Mahon}
\date{}

\begin{document}
\maketitle

\begin{quote}
\small
\textbf{2020 Mathematics Subject Classification.}
03C20 (primary); 28A60, 60A10 (secondary).

\smallskip
\textbf{Keywords.}
countable additivity, first-order axiomatizability, ultraproduct,
Boolean algebra, finitely additive measure, \L{}o\'s's theorem.
\end{quote}

\begin{abstract}
In the first-order language of Boolean algebras equipped with a normalized finitely
additive charge, no first-order theory characterizes the models in which the charge
is $\sigma$-additive.  In particular, countable additivity is not first-order
axiomatizable among finitely additive probability Boolean algebras.  The proof is a
single ultraproduct argument via \L{}o\'{s}'s theorem: Dirac masses on $\mathcal{P}(\N)$
are individually $\sigma$-additive, but the charge induced on the diagonal copy of
$\mathcal{P}(\N)$ in their ultraproduct is the $\{0,1\}$-valued ultrafilter charge,
which is purely finitely additive --- contradicting \L{}o\'{s}'s theorem if any
first-order axiomatization existed.
\end{abstract}

A property of structures is first-order axiomatizable in a language $L$ if there exists
a set of $L$-sentences whose models are exactly the structures satisfying it.
\L{}o\'{s}'s theorem gives the key diagnostic: if a property fails in an ultraproduct
of structures that individually satisfy it, no first-order axiomatization exists.

Countable additivity distinguishes $\sigma$-additive measures from finitely additive
charges.  That it cannot be expressed by a single first-order sentence is clear: the
condition quantifies over countable sequences, unavailable in first-order logic.  Its
non-first-order character is treated as background knowledge in the probability logic
literature~\cite{Hoover1978,Keisler1985,FaginHalpernMegiddo1990}.

The philosophical significance of this gap was articulated by
Howson~\cite{Howson2008}, who observed that de~Finetti's consistency criterion
(finite additivity) shares two structural features with first-order
logic --- absoluteness and compactness --- and that both fail for
$\sigma$-additivity.  He identified a missing ``completeness theorem telling us
that the rules of probability extend to finite but not countable additivity''
\cite[p.~17]{Howson2008}.  The ultraproduct argument below makes this
observation precise: the class of $\sigma$-additive structures is not
elementary, so no first-order coherence condition can enforce
$\sigma$-additivity.\footnote{A parallel non-axiomatizability result holds in the
pure Boolean algebra language (without charge predicates): the property of
admitting a strictly positive $\sigma$-additive probability --- characterized
algebraically by Kantorovich--Vulikh--Pinsker (Dedekind $\sigma$-completeness $+$
weak $(\sigma,\infty)$-distributivity $+$ chargeability;
see~\cite[391D]{Fremlin3}) --- is not first-order axiomatizable.  This follows
from Tarski's classification of atomless Boolean algebras: the Lebesgue measure
algebra satisfies all three KVP conditions, while the Cohen algebra
$\mathrm{RO}(2^{<\omega})$ is Dedekind complete and chargeable but fails weak
$(\sigma,\infty)$-distributivity; yet both are atomless and hence elementarily
equivalent.}

We record a short self-contained ultraproduct proof in the
Boolean-algebra-with-charge setting, showing non-axiomatizability by \emph{any}
set of sentences, not merely a single formula.

\medskip

Let $\Lmu$ be the one-sorted first-order language with:
\begin{itemize}
  \item Boolean operations $\wedge, \vee, \neg, 0, 1$;
  \item for each rational $q \in [0,1] \cap \mathbb{Q}$, unary relation symbols
        $M_{\leq q}(x)$ and $M_{\geq q}(x)$, intended to mean
        $\mu(x) \leq q$ and $\mu(x) \geq q$ respectively.
\end{itemize}

A structure for $\Lmu$ is a Boolean algebra $B$ together with a $[0,1]$-valued finitely
additive probability $\mu$, encoded through the truth of the predicates $M_{\leq q}$
and $M_{\geq q}$.  Using rational comparison predicates rather than a function symbol
into $[0,1]$ avoids introducing a second numeric sort.

\begin{definition}[First-order theory of finitely additive probabilities]
Let $\Tfa$ be the $\Lmu$-theory axiomatizing:
\begin{enumerate}[label=(\roman*)]
  \item $B$ is a Boolean algebra;
  \item normalization: $\mu(0) = 0$ and $\mu(1) = 1$;
  \item monotonicity: $x \leq y \Rightarrow \mu(x) \leq \mu(y)$;
  \item finite additivity: for all rationals $r, s$ with $r + s \leq 1$,
        \[
          x \wedge y = 0 \wedge M_{\leq r}(x) \wedge M_{\leq s}(y)
          \;\Rightarrow\; M_{\leq r+s}(x \vee y),
        \]
        and similarly for lower bounds.
\end{enumerate}
The theory $\Tfa$ may be supplemented with rational-cut coherence axioms
(monotonicity of the predicates in $q$, and $M_{\leq 1}(x)$, $M_{\geq 0}(x)$
for all $x$) to ensure that the predicates determine a unique value
$\mu(x) \in [0,1]$ for each element; the proof below does not depend on whether
these are included.
\end{definition}

\noindent
The class of Boolean algebras with normalized finitely additive probability is
first-order axiomatizable in $\Lmu$.  Countable additivity, by contrast, requires
quantification over an arbitrary countable sequence $(a_n)_{n \in \N}$ --- not
available in first-order logic --- and so is not a sentence of $\Lmu$.

\medskip

\begin{proposition}[$\sigma$-additivity is not first-order axiomatizable]
\label{prop:not-fo}
There is no first-order $\Lmu$-theory $T \supseteq \Tfa$ such that, for every
$(B, \mu) \models \Tfa$,
\[
  (B, \mu) \models T \quad\Longleftrightarrow\quad \mu \text{ is $\sigma$-additive}.
\]
In particular, countable additivity is not first-order axiomatizable among
finitely additive probability Boolean algebras.
\end{proposition}

\begin{proof}
Let $\delta_n$ be the Dirac probability at $n \in \N$ on the $\sigma$-algebra
$\mathcal{P}(\N)$.  Each $(\mathcal{P}(\N), \delta_n)$ is a model of $\Tfa$ in
which $\delta_n$ is $\sigma$-additive.  Suppose for contradiction that a theory
$T$ as described exists; then each $(\mathcal{P}(\N), \delta_n) \models T$.

Let $\mathcal{U}$ be a nonprincipal ultrafilter on $\N$.  By \L{}o\'{s}'s
theorem \cite{Los1955}, the ultraproduct
$\prod_\mathcal{U}(\mathcal{P}(\N),\delta_n)$ also models $T$.  The charge
induced on the diagonal copy of $\mathcal{P}(\N)$ is
\[
  \ell(A) = \lim_\mathcal{U} \delta_n(A) = \lim_\mathcal{U} \mathbf{1}_{n \in A}.
\]
This is the $\{0,1\}$-valued charge corresponding to $\mathcal{U}$: $\ell(A) = 1$
iff $A \in \mathcal{U}$.  Since $\mathcal{U}$ is nonprincipal, $\ell(\{n\}) = 0$
for every $n$, yet $\ell(\N) = 1$.  Thus the sets $A_k = \{k, k+1, \ldots\}$
satisfy $A_0 \supseteq A_1 \supseteq \cdots$ and $\bigcap_k A_k = \varnothing$,
yet $\ell(A_k) = 1$ for all $k$ (since $A_k \in \mathcal{U}$ for all $k$, as
$\mathcal{U}$ is nonprincipal).  Hence $\ell$ is not $\sigma$-additive.

If the ultraproduct's charge were $\sigma$-additive, its restriction to the
diagonal copy of $\mathcal{P}(\N)$ would be $\sigma$-additive; but this
restriction is $\ell$, which is not.  Hence the ultraproduct does not satisfy
$\sigma$-additivity, contradicting \L{}o\'{s}'s theorem applied to the
assumption that each $(\mathcal{P}(\N), \delta_n) \models T$.
\end{proof}

\begin{remark}
The ultraproduct $\prod_\mathcal{U}(\mathcal{P}(\N),\delta_n)$ is an
$\Lmu$-structure in its own right; $\ell$ is the charge induced on the diagonal
copy of $\mathcal{P}(\N)$ inside it.  The diagonal copy is isomorphic to
$\mathcal{P}(\N)$ as a Boolean algebra, and $\ell$ is the $\{0,1\}$-valued
ultrafilter charge on it: a purely finitely additive probability that is not
$\sigma$-additive.
\end{remark}

\medskip

The proof exhibits a purely finitely additive charge --- the ultrafilter charge on
$\mathcal{P}(\N)$ --- as a $\mathcal{U}$-ultralimit of $\sigma$-additive
probabilities.  This shows that the class of $\sigma$-additive structures is not
closed under ultraproducts, hence not first-order axiomatizable.  A natural
follow-up question is which purely finitely additive charges on a Boolean algebra
arise as pointwise ultralimits of $\sigma$-additive probabilities on the same
algebra; the answer depends sensitively on the algebra's structural properties.
We note only that on a general Boolean algebra $\sigma$-additive probabilities may
fail to exist at all: Gaifman~\cite{Gaifman1964} exhibits, in ZFC, an atomless
Boolean algebra carrying no strictly positive finitely additive (a fortiori no
$\sigma$-additive) measure.

\bibliographystyle{plain}
\bibliography{../references}

\end{document}
